\chapter{Analýza}
\label{ch:analyza}

\section{Obezita, nadváha a význam sledování stravy}
\label{sec:obezita}

Obezita a nadváha patří mezi klíčové determinanty morbidity a mortality v rozvinutých i rozvojových zemích.
Dlouhodobě zvýšená tělesná hmotnost je spojena se zvýšeným rizikem kardiovaskulárních onemocnění, diabetu 2.~typu, některých nádorových onemocnění a zhoršenou kvalitou života.
Epidemiologická data Eurostatu (zjišťování EU-SILC, referenční rok 2022) ukazují, že Česká republika patří v~rámci Evropské unie k~zemím s~vyšším podílem obézních dospělých~\cite{eurostatBmi2022}.
Konkrétní hodnoty byly uvedeny v~kapitole~\ref{ch:uvod}.

\subsection{Index tělesné hmotnosti a jeho interpretace}
\label{subsec:bmi}

Základním orientačním ukazatelem pro posouzení tělesné hmotnosti je index tělesné hmotnosti (Body Mass Index, BMI).
BMI se vypočítá jako podíl tělesné hmotnosti v~kilogramech a druhé mocniny tělesné výšky v~metrech:

\begin{equation}
	\mathrm{BMI} = \frac{m}{h^{2}},
	\label{eq:bmi}
\end{equation}

\noindent kde $m$ je tělesná hmotnost v~kilogramech a $h$ tělesná výška v~metrech.
Na základě hodnoty BMI se rozlišují jednotlivé kategorie výživového stavu definované Světovou zdravotnickou organizací, které jsou shrnuty v~tabulce~\ref{tab:bmi}.

\begin{table}[ht]
	\centering
	\caption{Hodnoty BMI a jejich interpretace.}
	\label{tab:bmi}
	\begin{tabular}{ll}
		\toprule
		\textbf{BMI (kg/m\textsuperscript{2})} & \textbf{Kategorie} \\
		\midrule
		pod 18,5 & podváha \\
		18,5--24,9 & normální hmotnost \\
		25,0--29,9 & nadváha \\
		30,0--34,9 & obezita I.~stupně \\
		35,0--39,9 & obezita II.~stupně \\
		40,0 a více & obezita III.~stupně \\
		\bottomrule
	\end{tabular}
\end{table}

BMI je pro populační studie a rychlé screeningové hodnocení velmi užitečné, protože umožňuje jednoduché srovnání mezi jedinci i populacemi.
Současně má však významná omezení.
Nezohledňuje rozložení tuku v~těle, podíl svalové hmoty ani individuální rozdíly v~tělesném složení.
U sportovců nebo lidí s~vysokým podílem svalové hmoty může proto docházet k~systematické klasifikaci do nadváhy či obezity, ačkoli jejich zdravotní riziko tomu neodpovídá~\cite{theodorakis2025humanEnergyBalance}.
BMI by tak měl být v~klinické i praktické práci vnímán jako orientační ukazatel, který je vhodné doplnit dalšími měřeními (obvod pasu, složení těla, laboratorní ukazatele).

\subsection{Životní styl, energetická bilance a rozvoj obezity}
\label{subsec:zivotniStyl}

Současné konceptuální modely chápou obezitu jako dlouhodobý důsledek pozitivní energetické bilance, kdy energetický příjem převyšuje energetický výdej.
Theodorakis a~Nikolaou popisují, že do energetické bilance vstupuje řada vzájemně provázaných proměnných, například bazální metabolismus, termický efekt potravy, spontánní pohyb, strukturovaná fyzická aktivita a adaptivní metabolické změny~\cite{theodorakis2025humanEnergyBalance}.
To znamená, že dvě osoby se stejným kalorickým příjmem a podobnou tělesnou hmotností mohou mít odlišný dlouhodobý vývoj hmotnosti v~závislosti na rozdílech v~energetickém výdeji a hormonální regulaci.

Drenowatz a~Greier zdůrazňují, že efektivní řízení tělesné hmotnosti vyžaduje současnou práci se stravou i fyzickou aktivitou.
Izolované intervence zaměřené pouze na dietní omezení nebo pouze na cvičení dosahují obvykle menšího a hůře udržitelného efektu než integrované přístupy, které systematicky pracují s~oběma složkami energetické bilance~\cite{drenowatz2025integrating}.
Z~hlediska návrhu digitálních nástrojů je proto důležité uvažovat tělesnou hmotnost v~širším kontextu životního stylu, nikoli jen jako mechanický výsledek krátkodobého kalorického deficitu.

\section{Self monitoring jako nástroj řízení hmotnosti}
\label{sec:selfMonitoring}

Jedním z~klíčových behaviorálních principů používaných v~programech redukce hmotnosti je self monitoring, tedy systematické sledování vlastního chování.
V~oblasti výživy se self monitoring nejčastěji projevuje jako pravidelný záznam konzumovaných jídel, sledování tělesné hmotnosti a zaznamenávání fyzické aktivity.
Cílem je zvýšit povědomí o~vlastních návycích a umožnit uživateli, aby své chování průběžně upravoval v~souladu s~doporučeními.

Systematické přehledy ukazují, že pravidelný self monitoring příjmu potravy a tělesné hmotnosti je spojen s~větším úbytkem hmotnosti a lepším udržením dosažených změn než situace, kdy je sledování nepravidelné nebo zcela chybí.
Zároveň se ukazuje, že self monitoring lze realizovat různými způsoby od papírových deníků až po digitální nástroje a mobilní aplikace, přičemž forma záznamu ovlivňuje jak přesnost dat, tak dlouhodobou udržitelnost celého procesu.
Detailní přehled vybraných mobilních aplikací pro sledování stravy a jejich charakteristik je proto věnován samostatné podkapitole~\ref{sec:existujici}~\cite{ang2021efficacy,raber2021systematicReview}.

Souhrnně lze říci, že obezita a nadváha představují významný zdravotní problém v~České republice i celosvětově a že self monitoring patří mezi hlavní behaviorální nástroje pro redukci a udržení tělesné hmotnosti.
Jeho dlouhodobá účinnost je však omezená tím, že pravidelné zaznamenávání jídel a souvisejících údajů je časově i kognitivně náročné, což vede u~řady uživatelů k~postupnému poklesu frekvence záznamů a k~oslabení efektu jinak dobře nastavených intervencí.
Tento kontext vytváří přímou motivaci pro hledání nástrojů, které sníží zátěž spojenou se sledováním stravy a tím zvýší pravděpodobnost, že uživatelé budou nástroj využívat konzistentně i v~delším časovém horizontu~\cite{pujia2025roleMobileApps,crane2023patterns,theodorakis2025humanEnergyBalance}.

\section{Energetická bilance a nutriční základy}
\label{sec:energetickaBilance}

Pochopení principu energetické bilance je zásadní pro interpretaci výstupů z~každé aplikace, která pracuje s~kalorickým příjmem a výdejem.
Následující podkapitoly definují klíčové pojmy: energetickou hodnotu makroživin, složky celkového energetického výdeje, prediktivní rovnice pro odhad bazálního metabolismu, samotný vztah mezi příjmem a výdejem a zdroje chyb v~odhadu obou veličin.

\subsection{Kalorie a energetická hodnota potravin}
\label{subsec:kalorie}

V~nutriční praxi se jako jednotka energie používá kilokalorie (kcal), případně kilojoule (kJ).
Jedna kilokalorie odpovídá množství energie potřebné ke zvýšení teploty jednoho kilogramu vody o~jeden stupeň Celsia.
V~běžné komunikaci se pojem „kalorie" používá téměř vždy ve významu kilokalorie.
Energetická hodnota potravin vyjadřuje množství energie, které organismus může z~dané potraviny získat.
Základní makroživiny se liší i svou energetickou hodnotou.
Sacharidy a bílkoviny poskytují přibližně 4~kcal na gram, zatímco tuky zhruba 9~kcal na gram, tedy více než dvojnásobek~\cite{theodorakis2025humanEnergyBalance}.

Pro udržení tělesné hmotnosti v~rovnováze musí být dlouhodobě energetický příjem z~potravin v~rovnováze s~energetickým výdejem.
Systematické odchylky v~podobě chronického kalorického deficitu nebo nadbytku vedou k~postupné změně tělesné hmotnosti~\cite{theodorakis2025humanEnergyBalance}.

\subsection{Složky energetického výdeje}
\label{subsec:slozkyVydeje}

Celkový denní energetický výdej (TEE) se obvykle rozděluje na několik složek:

\begin{itemize}
	\item bazální metabolismus (BMR),
	\item termický efekt potravy (TEF),
	\item energetický výdej na fyzickou aktivitu, který zahrnuje jak strukturované cvičení, tak běžný denní pohyb.
\end{itemize}

Theodorakis a~Nikolaou uvádějí, že u~většiny dospělé populace tvoří bazální metabolismus přibližně 60 až 75~\% celkového energetického výdeje.
Zbývající část připadá na termický efekt potravy a fyzickou aktivitu, přičemž jejich relativní význam se může výrazně lišit v~závislosti na životním stylu~\cite{theodorakis2025humanEnergyBalance}.

Bazální metabolismus představuje množství energie, které organismus potřebuje v~klidových podmínkách pro udržení základních životních funkcí, jako je činnost srdce a plic, udržování tělesné teploty a syntéza tkání.
Je ovlivněn především:

\begin{itemize}
	\item množstvím a složením tělesné hmoty (zejména podílem beztukové hmoty),
	\item věkem,
	\item pohlavím,
	\item genetickými faktory a hormonálním stavem.
\end{itemize}

\subsection{Prediktivní rovnice pro odhad bazálního metabolismu}
\label{subsec:bmrRovnice}

Přímé měření energetického výdeje pomocí nepřímé kalorimetrie není v~běžné praxi ani v~uživatelských aplikacích reálně dostupné.
Místo něj se využívají empirické prediktivní rovnice odhadující BMR ze snadno měřitelných parametrů (hmotnost, výška, věk, pohlaví).
Historicky nejrozšířenější je Harris-Benedictova rovnice z~roku 1919, novější Mifflin-St~Jeorova rovnice z~roku 1990 však bývá řadou novějších studií hodnocena jako přesnější pro běžnou dospělou populaci~\cite{frankenfield2005comparison} a má tvar:

\begin{align}
	\mathrm{BMR}_{\text{muži}} &= (10 \cdot m) + (6{,}25 \cdot h) - (5 \cdot a) + 5,   \label{eq:msj-muzi} \\
	\mathrm{BMR}_{\text{ženy}} &= (10 \cdot m) + (6{,}25 \cdot h) - (5 \cdot a) - 161, \label{eq:msj-zeny}
\end{align}

\noindent kde $m$ je hmotnost v~kilogramech, $h$ výška v~centimetrech a $a$ věk v~letech.
Systematické přehledy zároveň ukazují, že přesnost prediktivních rovnic se liší mezi populacemi a u~některých skupin (zejména sportovci a lidé s~nadprůměrným podílem svalové hmoty) mohou běžně používané rovnice BMR systematicky podhodnocovat nebo nadhodnocovat~\cite{frankenfield2005comparison}.
V~uživatelských aplikacích je proto použití těchto rovnic praktickým kompromisem mezi dostupností dat a požadavkem na rychlý odhad energetického výdeje, doprovázeným srozumitelnou komunikací nejistoty výsledku směrem k~uživateli.

\subsection{Energetická bilance a energetická nerovnováha}
\label{subsec:nerovnovaha}

Energetickou bilanci lze zjednodušeně vyjádřit rovnicí:

\begin{equation}
	\Delta E_{\text{zásoby}} = E_{\text{příjem}} - E_{\text{výdej}}.
	\label{eq:bilance}
\end{equation}

Pokud je energetický příjem dlouhodobě nižší než energetický výdej, organismus čerpá chybějící energii z~uložených zásob, především z~tukové tkáně.
Tento stav se označuje jako kalorický deficit a vede k~postupnému snižování tělesné hmotnosti.
Příliš výrazný a dlouhodobý deficit však může kromě tukové tkáně ohrožovat i svalovou hmotu a vést k~nedostatku živin, únavě a zhoršení zdravotního stavu~\cite{theodorakis2025humanEnergyBalance,drenowatz2025integrating}.

Naopak při chronickém kalorickém nadbytku převyšuje energetický příjem výdej.
Přebytečná energie se ukládá ve formě tukových zásob a dochází k~postupnému zvyšování tělesné hmotnosti a rizika metabolických onemocnění.
Drenowatz a~Greier zdůrazňují, že významnou roli hraje nejen absolutní množství energie, ale také struktura stravy, rozložení příjmu během dne a míra fyzické aktivity.
Účinné strategie řízení hmotnosti proto kombinují dietní úpravy s~podporou pohybové aktivity~\cite{drenowatz2025integrating}.

\subsection{Zdroje chyb v~odhadu příjmu a výdeje energie}
\label{subsec:zdrojeChyb}

Praktické využití energetické bilance v~každodenním životě naráží na zásadní problém, kterým je nepřesnost odhadů jak na straně příjmu, tak na straně výdeje.
Uživatelé běžně podhodnocují množství a energetickou hodnotu zkonzumovaných potravin, zapomínají zaznamenávat některá jídla nebo špatně odhadují velikost porcí.
Li a~kol. například prokázali, že výstupy nutričních aplikací se mohou významně lišit od referenčních metod, a to jak při manuálním zapisování, tak při využití rozpoznávání obrazu~\cite{li2024evaluating}.

Na straně energetického výdeje může docházet k~nadhodnocování spálené energie, zejména pokud uživatel spoléhá na obecné tabulky nebo méně přesná zařízení.
Theodorakis a~Nikolaou upozorňují, že organismus navíc na energetický deficit adaptivně reaguje snížením bazálního metabolismu a spontánního pohybu, což dále komplikuje přesnou predikci změn tělesné hmotnosti~\cite{theodorakis2025humanEnergyBalance}.

Pro návrh digitálního nástroje to znamená, že:

\begin{itemize}
	\item nelze spoléhat na absolutní přesnost jednoho konkrétního čísla energetického příjmu nebo výdeje,
	\item je vhodné pracovat s~konceptem odhadu s~nejistotou,
	\item design rozhraní by měl uživatele motivovat ke konzistentnímu zapisování a zároveň usilovat o~minimalizaci zátěže.
\end{itemize}

Právě tato snaha o~snížení uživatelské zátěže při zachování dostatečné přesnosti odhadu energetické bilance je jedním z~hlavních důvodů, proč tato práce zkoumá možnosti využití umělé inteligence a rozpoznávání potravin z~fotografie.

\section{Makroživiny a mikroživiny ve výživě}
\label{sec:zivoiny}

Při hodnocení kvality stravy je vedle celkového energetického příjmu zásadní také její složení.
To určují především makroživiny, které dodávají energii a stavební materiál, a mikroživiny, jež umožňují správný průběh metabolických a regulačních procesů.
Přehledové studie ukazují, že jak poměr makroživin, tak dostatečný příjem klíčových mikroživin mají významný vliv na tělesné zdraví, riziko chronických onemocnění i subjektivní pohodu~\cite{venn2020macronutrients,farag2021dietaryMicronutrients}.

\subsection{Makroživiny}
\label{subsec:makrozivny}

Makroživiny tvoří hlavní zdroj energie ve stravě a zahrnují sacharidy, tuky a bílkoviny.
Sacharidy (jednoduché cukry, škrob, vláknina) jsou primárním zdrojem energie pro většinu buněk, přičemž strava s~vyšším zastoupením komplexních sacharidů a vlákniny je spojována s~lepší kontrolou glykemie a vyšším pocitem sytosti.
Tuky představují koncentrovaný zdroj energie a jsou nezbytné pro vstřebávání lipofilních vitamínů.
Významnou roli hraje jejich kvalitativní složení, neboť vyšší podíl nenasycených mastných kyselin je spojován s~příznivým vlivem na kardiovaskulární zdraví~\cite{venn2020macronutrients}.
Bílkoviny jsou základním stavebním materiálem tkání a hormonů, přičemž jejich dostatečný příjem je klíčový zejména pro udržení svalové hmoty při redukci tělesné hmotnosti~\cite{venn2020macronutrients}.

Z~praktického hlediska je vedle celkového energetického příjmu důležitá i kvalita jednotlivých makroživin a dlouhodobá udržitelnost stravovacího režimu, protože různé dietní přístupy mohou vést ke srovnatelným změnám hmotnosti, pokud je zachován celkový energetický deficit~\cite{theodorakis2025humanEnergyBalance,drenowatz2025integrating}.

\subsection{Mikroživiny}
\label{subsec:mikrozivny}

Mikroživiny (vitaminy a minerální látky) na rozdíl od makroživin nepřinášejí významné množství energie, ale podílejí se na enzymatických a regulačních procesech v~organismu a jejich nedostatek může ovlivnit imunitu, kognitivní funkce i riziko chronických onemocnění~\cite{farag2021dietaryMicronutrients}.
V~aktuální verzi navrhované aplikace se sledování mikroživin neuvažuje, neboť přesahuje rámec primárního cíle, kterým je snížení bariéry kalorického záznamu.
Jejich začlenění je však relevantním kandidátem pro budoucí rozšíření aplikace o~kvalitativní hodnocení stravy, zejména v~kontextu fenoménu \uv{skryté malnutrice}, kdy vysoký energetický příjem může být doprovázen nedostatečným příjmem některých mikroživin~\cite{farag2021dietaryMicronutrients}.

\section{Existující mobilní aplikace pro sledování stravy}
\label{sec:existujici}

Mobilní aplikace pro sledování stravy a tělesné hmotnosti patří v~současnosti k~nejčastěji používaným nástrojům digitálního zdraví a přehledové studie potvrzují jejich přínos pro lepší kontrolu stravovacích návyků a mírnou redukci hmotnosti~\cite{nogueira2024mobileApplications}.
Systematické průzkumy hlavních mobilních obchodů identifikují v~této kategorii stovky vzájemně konkurujících aplikací, například Samad a~kol. nalezli 473 aplikací souvisejících se sledováním příjmu potravy, z~nichž 80 podrobně analyzovali~\cite{samad2022smartphoneApps,abeltino2025digitalApplications}.

Přehledy výživových aplikací ukazují, že většina z~nich staví na velmi podobném jádru a liší se spíše v~uživatelském rozhraní a doplňkových funkcích než v~samotné logice sledování stravy~\cite{nogueira2024mobileApplications,abeltino2025digitalApplications}.
Společné jádro typicky zahrnuje potravinový deník, databázi potravin s~fulltextovým vyhledáváním, ukládání vlastních receptů a oblíbených jídel, nastavení cílového energetického příjmu a základní vizualizaci výsledků formou grafů a statistik~\cite{gioia2023mobileApps}.
Z~hlediska zadávání dat většina komerčních aplikací nadále spoléhá na manuální zapisování doplněné o~částečně automatizované zkratky, zejména skenování čárových kódů, šablony jídel s~možností duplikace dříve zapsaných záznamů a připomenutí ve formě notifikací~\cite{abeltino2025digitalApplications}.
Empirickou oporu pro popis silných a slabých stránek této kategorie aplikací poskytuje analýza Zečević a~kol., kteří metodou \emph{topic modelingu} zpracovali 72\,084 recenzí z~Google Play Store pro patnáct nejpoužívanějších aplikací pro sledování stravy a identifikovali jedenáct opakujících se témat~\cite{zecevic2021userPerspectives}.
Mezi nejhůře hodnocené patří netransparentní účtování předplatného (průměrné hodnocení 2,32 z~5) a technická stabilita (2,85 z~5), zatímco kladně jsou vnímány jednoduchost ovládání, samotná funkce zápisu kalorií a integrace čtečky čárových kódů~\cite{zecevic2021userPerspectives}.

Systematické studie zároveň upozorňují, že ruční zapisování je časově náročné a v~praxi často vede k~nekompletním či zkresleným záznamům.
Li a~kol. ukazují, že výstupy nutričních aplikací se mohou od referenčních metod významně lišit jak u~manuálního zapisování, tak u~řešení s~podporou rozpoznání obrazu, přičemž energetický příjem bývá podle typu stravy nadhodnocen nebo podhodnocen o~stovky kilokalorií denně~\cite{li2024evaluating}.
V~reakci na tato omezení se prosazuje využití umělé inteligence a počítačového vidění, avšak Samad a~kol. ve svém vzorku 80 aplikací identifikovali pouze jedinou, která dokázala z~fotografie automaticky určit druh jídla i velikost porce~\cite{samad2022smartphoneApps}.
Aburub a~kol. shrnují, že plně automatické rozpoznávání talíře zůstává výjimkou a typicky vyžaduje potvrzení a korekci uživatele, zatímco skenování čárových kódů se mezitím stalo běžným standardem~\cite{aburub2025foodScanning}.

Chotwanvirat a~kol. ve své rešerši 84 prací publikovaných do roku 2021 dokumentují, že po roce 2015 vystřídaly ruční extrakci příznaků konvoluční neuronové sítě překračující 80~\% přesnosti na datasetech FOOD-101 a UEC FOOD-256~\cite{chotwanvirat2024scoping}.
Tatáž rešerše upozorňuje, že většina takto trénovaných systémů zůstává omezena na curated datasety, hůře se vyrovnává s~reálnými podmínkami snímání a postrádá vysvětlitelnost svých rozhodnutí, což ztěžuje jejich přijetí v~klinické praxi~\cite{chotwanvirat2024scoping}.
Liu a~kol. tento závěr kvantifikují i na novější transformerové architektuře CBiAFormer, jejíž přesnost klesá z~92,40~\% na západním datasetu ETHZ Food-101 na 72,92~\% na smíšeném ISIA Food-200, což ilustruje obecnější obtíž specializovaných modelů zobecnit napříč kuchyněmi~\cite{liu2025deepLearning}.

Po roce 2022 nastupuje v~oboru třetí vlna v~podobě multimodálních velkých jazykových modelů, které dokáží v~jediném volání zpracovat obraz, text i hlasový přepis a vrátit strukturovaný odhad nutričních hodnot bez nutnosti vlastního trénování na potravinovém datasetu~\cite{wang2026visionLanguageFood,fridolfsson2025llmNutrition}.
Recentní studie a preprinty z~let 2025 a 2026 ukazují, že tyto modely při vhodném prompt engineeringu dosahují srovnatelné nebo lepší přesnosti odhadu kalorické hodnoty než dříve nasazované specializované sítě~\cite{hua2025nutribench,coburn2025multimodalNutrition}.
Tento posun zároveň přesouvá těžiště přesnosti z~architektury a tréninkových dat směrem k~designu promptu a postprocessingu odpovědi modelu, čemuž je v~rámci této práce věnována experimentální validace popsaná v~kapitole~\ref{ch:testovani}~\cite{ase2025structuredPrompts}.

Pro tuto práci jsou reprezentativní čtyři aplikace, které lze rozdělit do dvou typologicky odlišných skupin.
První skupinu tvoří tradiční aplikace s~manuálním potravinovým deníkem, v~němž umělá inteligence funguje jako asistent rychlejšího zápisu, zastoupené globálním MyFitnessPal a československými Kalorickými Tabulkami (Dine4Fit).
Druhou skupinu tvoří novější aplikace stavící primárně na rozpoznávání jídla z~fotografie, ilustrované americkým SnapCalorie a novějším komerčním zástupcem této vlny CalZen.
Toto rozdělení odráží rozdíl mezi etablovaným ekosystémem deníkových aplikací a komerčním projevem třetí vlny popsané výše, jejíž jádro tvoří multimodální modely nasazené po roce 2022~\cite{wang2026visionLanguageFood,fridolfsson2025llmNutrition}.
Následující podsekce obě skupiny stručně představují a vymezují jejich silné a slabé stránky ve vztahu k~cíli této diplomové práce.

\subsection{Tradiční aplikace s~manuálním zápisem}
\label{subsec:tradicniAplikace}

Tato skupina představuje dlouhodobě dominantní jádro trhu, na němž stojí drtivá většina z~více než čtyř stovek aplikací identifikovaných systematickými průzkumy hlavních mobilních obchodů~\cite{samad2022smartphoneApps,abeltino2025digitalApplications}.
Jejím společným rysem je strukturovaný potravinový deník navázaný na rozsáhlou databázi potravin a zápis založený především na manuálním vyhledávání, doplněný o~částečně automatizované zkratky typu skenování čárových kódů, šablon jídel a foto rozpoznávání~\cite{abeltino2025digitalApplications,gioia2023mobileApps}.
Pro tuto skupinu je charakteristické, že umělá inteligence vstupuje až jako sekundární funkce nad existujícím deníkovým rozhraním, nikoli jako primární vstupní kanál~\cite{aburub2025foodScanning}.

\subsubsection{MyFitnessPal}
\label{subsubsec:myfitnesspal}

MyFitnessPal patří mezi nejznámější a nejrozšířenější aplikace pro sledování stravy a tělesné hmotnosti na světě, je dostupný pro Android, iOS i webové rozhraní a dlouhodobě se pohybuje na předních příčkách kategorií Health \& Fitness v~hlavních mobilních obchodech~\cite{myFitnessPal2025}.
Základem aplikace je potravinový deník navázaný na rozsáhlou databázi potravin obsahující miliony položek včetně značkových výrobků a restaurací.
Uživatel zapisuje jídla vyhledáváním v~databázi, výběrem z~historie, skenováním čárového kódu nebo zadáním vlastních receptů.
Aplikace průběžně počítá denní energetický příjem a rozložení makroživin, podporuje záznam fyzických aktivit, sledování hmotnosti a integraci s~externími zařízeními (Apple Health, fitness náramky).
Pokročilé funkce, jako je sledování mikronutrientů a detailnější nutriční přehledy, jsou součástí placeného předplatného Premium~\cite{myFitnessPal2025}.

MyFitnessPal je jedním z~prvních velkých komerčních produktů, které nasadily ve větším měřítku foto rozpoznávání jídel.
Funkce Meal Scan umožňuje uživateli namířit fotoaparát na talíř, přičemž aplikace pomocí počítačového vidění navrhne odpovídající položky a velikosti porcí z~interní databáze, které uživatel potvrdí nebo upraví~\cite{myFitnessPal2025}.
Klíčovými přednostmi z~hlediska této práce jsou rozsáhlá globální databáze, kombinace více vstupních metod a dlouhodobá přítomnost na trhu doložená nezávislými studiemi~\cite{samad2022smartphoneApps}.
Mezi omezeními naopak figuruje proměnlivá kvalita uživatelsky přidaných položek vedoucí k~chybám v~nutričních údajích, nutnost ruční kontroly výsledků foto rozpoznávání~\cite{li2024evaluating} a nedostatečné pokrytí lokálních středoevropských produktů.
MyFitnessPal je tak referenčním příkladem robustního ekosystému, který umělou inteligenci využívá primárně jako doplněk manuálního zápisu, nikoli jako jeho plnou náhradu.

\subsubsection{Kalorické Tabulky}
\label{subsubsec:kalorickeTabulky}

Kalorické Tabulky, vyvíjené společností Dine4Fit~a.s., představují nejpoužívanější aplikaci pro sledování stravy v~českém a slovenském prostředí.
Podle oficiálních údajů firmy ji měsíčně používá více než jeden milion uživatelů~\cite{dine4fit2025}.
Aplikace je dostupná jako web i mobilní aplikace pro Android a iOS a podporuje vybrané chytré hodinky (Garmin, Apple Watch).
Její základní princip se podobá globálním nástrojům: potravinový deník navázaný na databázi výrazně přizpůsobenou regionálnímu trhu, sledování pitného režimu a hmotnosti, záznam fyzických aktivit a kalorická kalkulačka odvozující doporučený příjem z~tělesných parametrů a cílů uživatele~\cite{dine4fit2025}.
Stejně jako MyFitnessPal funguje aplikace ve freemium modelu, kde rozšířené analýzy, mikronutrienty a pokročilé statistiky jsou součástí placeného Premium~\cite{dine4fit2025}.

Novější verze Kalorických Tabulek integrují funkci rozpoznávání jídel z~fotografie pomocí umělé inteligence: uživatel jídlo vyfotí a aplikace navrhne odpovídající potraviny a přibližnou gramáž~\cite{dine4fit2025}.
Stejně jako u~MyFitnessPal nejde o~plně autonomní systém, nezávislá kvantitativní studie přesnosti této konkrétní funkce v~recenzované literatuře zatím chybí, což odpovídá obecnějšímu trendu opožděné akademické validace komerčních AI funkcí~\cite{samad2022smartphoneApps,li2024evaluating,aburub2025foodScanning}.
Hlavními přednostmi v~českém kontextu jsou silná lokalizace databáze, dostupnost webového i mobilního rozhraní a velká uživatelská základna přispívající ke komunitním receptům~\cite{dine4fit2025}.
Mezi limity patří časová náročnost při práci se složitějšími recepty, riziko nepřesností v~komunitně rozšiřovaných databázích~\cite{li2024evaluating} a stále spíše doplňková role AI rozpoznávání vůči manuálnímu zápisu.

\subsection{Aplikace stavící primárně na AI rozpoznávání}
\label{subsec:aiFirstAplikace}

Skupina se vnitřně dělí na dva přístupy, které jsou v~následujících podkapitolách reprezentovány aplikacemi SnapCalorie a CalZen, totiž vlastní specializovaný model trénovaný na uceleném potravinovém datasetu a generický multimodální model nasazený přes API~\cite{coburn2025multimodalNutrition}.

\subsubsection{SnapCalorie}
\label{subsubsec:snapcalorie}

Aplikace SnapCalorie je vyvíjena americkým týmem, jehož zakladatelé pocházejí z~výzkumu počítačového vidění ve společnosti Google a podle veřejných materiálů firmy se podíleli na vzniku produktů Google Lens a Cloud Vision API~\cite{snapcalorie2025}.
Technologickým jádrem je vlastní specializovaný model trénovaný na proprietárním datasetu zhruba 5\,000 jídel, jejichž ingredience byly před fotografováním na robotickém rigu zváženy na laboratorní váze~\cite{snapcalorie2025}.
Cílem této koncepce je pokrýt obtížné případy typu polévek, mixovaných pokrmů a omáček, na kterých generické datasety selhávají~\cite{snapcalorie2025,chotwanvirat2024scoping}.
Doprovodnou výbavou je databáze přibližně 500\,000 položek validovaná oproti referenčnímu zdroji USDA, skener nutričních štítků a integrace s~platformou Apple HealthKit, přičemž každý odhad lze ručně upravit a model je deklarován jako učící se z~uživatelských korekcí~\cite{snapcalorie2025}.

Sebereportovaná průměrná chyba odhadu energetické hodnoty činí podle vyjádření výrobce přibližně 15~\%, což je v~publikovaných materiálech srovnáváno s~tolerancí 20~\% povolenou na nutričních štítcích a s~vyššími chybovostmi nutričních specialistů a běžných uživatelů konkurenčních aplikací~\cite{snapcalorie2025}.
Tyto údaje pocházejí výhradně od provozovatele a v~recenzované literatuře k~nim zatím chybí nezávislá validační studie, což odpovídá obecnějšímu jevu opožděné akademické verifikace komerčních AI funkcí~\cite{samad2022smartphoneApps,li2024evaluating,aburub2025foodScanning}.
Z~hlediska této práce je SnapCalorie referenčním příkladem \emph{AI-first} přístupu se specializovaným vlastním modelem, jehož hlavními limity jsou převážně západní složení trénovacího datasetu, slabá lokalizace mimo Spojené státy a obtížnější přenositelnost do prostředí české a středoevropské kuchyně, kterou kvantifikuje pokles přesnosti specializovaných sítí napříč regionálními datasety~\cite{liu2025deepLearning}.

\subsubsection{CalZen}
\label{subsubsec:calzen}

Aplikace CalZen reprezentuje generickou vlnu AI-first nástrojů uvedených v~letech 2024 a 2025, které slibují odhad kalorické a makroživinové hodnoty z~fotografie pokrmu řádově do několika sekund a podle vyjádření provozovatele rozpoznají domácí jídla, restaurační pokrmy i balené potraviny~\cite{calzen2025}.
Provozovatel deklaruje řádově tři miliony uživatelů a průměrné hodnocení 4,8~z~5 v~hlavních mobilních obchodech, jedná se však o~marketingové údaje výrobce bez nezávislé verifikace~\cite{calzen2025}.
Doprovodné funkce zahrnují krokoměr a sledování pitného režimu a aplikace funguje ve freemium modelu, v~němž jsou pokročilé makroanalýzy a integrace dostupné v~placeném tarifu~\cite{calzen2025}.

Na rozdíl od SnapCalorie firma technické řešení blíže nedokumentuje a nezveřejňuje, jaký model rozpoznávání používá ani jakou metodikou je měřena přesnost odhadu~\cite{calzen2025}.
Krátký inferenční čas a generická šíře pokrývající různorodé kuchyně, restaurace i obalová balení jsou konzistentní s~nasazením multimodálních velkých jazykových modelů přes API, jak je dokumentuje recentní literatura~\cite{wang2026visionLanguageFood,fridolfsson2025llmNutrition,coburn2025multimodalNutrition}.
Pro tuto práci je CalZen typovým představitelem druhého přístupu v~rámci AI-first skupiny, tedy generického multimodálního modelu nasazeného bez vlastního trénování, jehož silnou stránkou je rychlost vývoje a šíře pokrytí, slabinami pak nízká technická transparentnost, absence veřejné validační studie a slabé pokrytí české kuchyně srovnatelné s~limity SnapCalorie~\cite{li2024evaluating,aburub2025foodScanning}.

\subsection{Shrnutí aplikací na trhu}
\label{subsec:shrnutiTrhu}

Čtyři představené aplikace ilustrují dvě paralelní vlny současného trhu se sledováním stravy, totiž tradiční deníkové nástroje s~rozsáhlými databázemi a milionovými uživatelskými základnami na straně jedné a AI-first aplikace stavící foto rozpoznávání jako primární vstup na straně druhé~\cite{abeltino2025digitalApplications,wang2026visionLanguageFood}.
Druhá skupina není inovací posledních měsíců, ale několikaletým komerčním trendem, který v~návaznosti na multimodální modely uvedené po roce 2022 už generuje uživatelské základny v~řádu milionů, jak dokládají sebereportované údaje aplikací typu SnapCalorie a CalZen~\cite{snapcalorie2025,calzen2025,fridolfsson2025llmNutrition}.
Společným rysem obou skupin však zůstává, že nezávislá akademická validace přesnosti komerčních AI funkcí v~recenzované literatuře nadále chybí, a to jak u~doplňkového foto rozpoznávání v~tradičních aplikacích typu MyFitnessPal Meal Scan, tak u~plně AI-first řešení~\cite{samad2022smartphoneApps,li2024evaluating,aburub2025foodScanning}.

Pozici této diplomové práce v~uvedeném kontextu nevymezuje ambice překonat komerční AI-first nástroje v~technické přesnosti rozpoznávání, ke které by autorský tým bez zdrojů firem typu SnapCalorie neměl reálné prostředky.
Diferenciace je postavena na třech osách, které ani jedna z~představených aplikací nepokrývá současně, totiž na lokalizované databázi potravin postavené na evropském zdroji Open Food Facts s~důrazem na regionální produkty, na transparentním zobrazení nejistoty AI výstupu uživateli a na integraci s~evropským zdravotním ekosystémem prostřednictvím Apple Health a Health Connect~\cite{dine4fit2025,liu2025deepLearning}.
Tato kombinace odpovídá identifikovaným slabinám trhu, na nichž zaostává jak globální MyFitnessPal v~lokalizaci, tak SnapCalorie a CalZen v~pokrytí české kuchyně i v~transparentnosti komunikace omezení modelu~\cite{zecevic2021userPerspectives,liu2025deepLearning}.

Z~uvedeného přehledu vyplývají pro návrh aplikace v~této diplomové práci tři hlavní závěry, které jsou konzistentní i s~kvantitativní analýzou nejčastějších stížností uživatelů ve studii Zečević a~kol.~\cite{zecevic2021userPerspectives}.
Za prvé, inovace by neměla cílit na samotnou existenci AI funkce, která je již na trhu komoditou, ale na konkrétní slabá místa stávajících řešení, zejména na uživatelskou zátěž zápisu, kvalitu lokální databáze a transparentní komunikaci nejistoty výstupu vůči uživateli~\cite{zecevic2021userPerspectives,li2024evaluating}.
Za druhé, AI-first přístup je dostupný i bez vlastního trénovaného modelu skrze multimodální velké jazykové modely nasazené přes API, což otevírá prostor pro rychlou implementaci v~rámci diplomové práce~\cite{wang2026visionLanguageFood}.
Za třetí, lokální kontext, podpora české a středoevropské kuchyně a kvalitní lokalizovaná databáze zůstávají diferenciátorem i vůči AI-first konkurenci typu SnapCalorie a CalZen, která je primárně orientovaná na západní a anglosaský trh~\cite{liu2025deepLearning,dine4fit2025}.

\section{Uživatelská studie}
\label{sec:uzivatelskaStudie}

Uživatelská studie představuje soubor metod, jejichž cílem je porozumět potřebám, motivacím a kontextu cílových uživatelů a tyto poznatky následně převést do návrhu systému.
V~oblasti mobilního zdraví je uživatelsky orientovaný přístup považován za důležitý, protože ovlivňuje použitelnost, důvěru, dlouhodobou adherenci a tím i reálný dopad aplikace~\cite{cornet2020untoldStories}.

V~této práci byla uživatelská studie použita jako hlavní vstup pro návrh aplikace pro rozpoznávání potravin pomocí umělé inteligence a pro sledování kalorického příjmu.
Cílem bylo identifikovat situace, ve kterých je ruční zapisování pro uživatele překážkou, jaká míra nepřesnosti je akceptovatelná, jak uživatelé vnímají selhání umělé inteligence a jaké podpůrné funkce zvyšují pravděpodobnost dlouhodobého používání.

Účastníci uživatelské studie jsou v~textu označováni \texttt{R1}--\texttt{R4} (Respondent rozhovorů) a tvoří odlišnou skupinu od participantů \texttt{P1}--\texttt{P3} v~kapitole~\ref{ch:testovani}, která referuje samostatné uživatelské testování implementované aplikace.

\subsection{Metodika}
\label{subsec:metodika}

Jako metoda sběru dat byly zvoleny strukturované rozhovory s~uživateli z~cílové skupiny.
Rozhovory byly vedeny tak, aby pokryly současnou praxi zapisování jídel, očekávání od rozpoznávání pomocí umělé inteligence a také preference v~oblasti přehledů, motivace, práce s~daty a integrací.
Tento postup odpovídá doporučením v~literatuře, kde jsou kvalitativní techniky typu rozhovorů nebo skupinových diskusí uváděny jako vhodný zdroj požadavků pro návrh aplikací v~oblasti zdraví~\cite{molina2020proposalUserCentered}.

Rozhovory byly zpracovány do samostatných reportů a dále analyzovány tematickým způsobem.
V~této kapitole jsou výsledky shrnuty do hlavních témat, která následně tvoří vstup pro formalizaci scénářů případů užití a sady funkčních a nefunkčních požadavků v~kapitole~\ref{ch:navrh}.
Reporty v~plném znění jsou součástí příloh této práce.

\subsection{Charakteristika respondentů}
\label{subsec:participanti}

Pro rozhovory byli vybráni čtyři respondenti, kteří reprezentují rozdílné cíle a kontext používání.
V~textu jsou uváděni anonymizovaně jako R1 až R4.

\begin{description}
	\item[R1] muž ve věku 27~let, jehož cílem je nabírání svalové hmoty.
		Cvičí přibližně třikrát týdně a preferuje sledování makroživin.
		V~minulosti aplikaci pro zapisování využíval, ale přestal, když dosáhl cíle.
		Je ochoten tolerovat určitou míru nepřesnosti, pokud mu aplikace ušetří práci.
	\item[R2] muž ve věku 24~let, intenzivně cvičí a cíleně udržuje kalorický nadbytek.
		U~zapisování klade důraz na přesnost a důvěryhodnost, preferuje rychlé metody zadávání u~balených potravin, typicky čárové kódy.
		U~jídel z~restaurace akceptuje typicky vyšší nepřesnost při zápisu než u~jiných pokrmů.
	\item[R3] žena ve věku 54~let, dlouhodobě se snaží zhubnout a zohledňuje zdravotní souvislosti.
		Zkoušela existující aplikace, ale odradila ji časová náročnost, časté chybění položek v~databázi a nutnost odhadovat gramáž, zejména u~domácího vaření.
	\item[R4] žena ve věku 28~let, která dříve své jídlo pravidelně vážila, nyní častěji zapisuje odhadem a upřednostňuje jednoduché předvolby porcí.
		Má jasnou preferenci, aby fotografie byla primárním vstupem, ale současně požaduje kvalitní záložní cestu při selhání.
		Zásadní je pro ni transparentnost umělé inteligence a rozlišení běžných limitů modelu od technické chyby aplikace.
\end{description}

\subsection{Výsledky rozhovorů}
\label{subsec:vysledkyRozhovoru}

\paragraph{Rychlost zápisu a dlouhodobá udržitelnost.}
Opakovaným tématem napříč rozhovory je, že dlouhodobé používání je přímo podmíněno rychlostí a jednoduchostí.
R1 zmiňuje cílový stav rychlého zápisu jednoduché položky přibližně do šesti kliknutí bez dodatečné editace.
R2 uvádí podobnou hranici pro rychlé položky.
R4 toleruje pomalejší zpracování zejména na začátku, ale očekává postupné zrychlení a udržení zápisu v~přiměřeném čase.
Zcela klíčové je to pro R3, která pomalost procesu zápisu označila za hlavní důvod, proč v~dosavadních pokusech přestávala aplikaci používat ještě před dosažením svých zdravotních cílů.

\paragraph{Akceptovatelná nepřesnost a důvěra v~umělou inteligenci.}
Respondenti se liší v~toleranci nepřesnosti, což je klíčové pro definici cílových metrik přesnosti.
R1 je ochoten tolerovat přibližnou odchylku do dvaceti procent, pokud mu aplikace výrazně zjednoduší proces.
Naopak R2 požaduje nepřesnost pouze v~řádu jednotek procent a bez důvěryhodného výsledku považuje funkci za málo užitečnou.
R3 chápe omezení vysoké přesnosti, nicméně požaduje, aby výstup byl srozumitelný a kontrolovatelný, zejména u~domácích jídel.

\paragraph{Selhání rozpoznání.}
Pro praktickou použitelnost funkce rozpoznávání je kritické eliminovat situace, kdy se uživatel ocitne ve slepé uličce.
R4 požaduje, aby při selhání fotografie bylo možné plynule přejít na textový popis a aby bylo možné rozpoznání znovu vyvolat z~rozpracovaného záznamu bez zakládání nového záznamu.
R1 zároveň navrhuje možnost doplnit k~fotografii krátkou doplňující informaci pro lepší výsledek.
Důležitým aspektem je také rozlišení, zda jde o~limit rozpoznání, nebo o~technickou chybu aplikace, a nabídnutí vhodného dalšího kroku.

\paragraph{Oblíbená jídla, historie a zkratky.}
Všichni respondenti popisují situace, kdy jedí opakovaně podobná jídla, případně používají stejné ingredience.
R2 označuje oblíbené položky, historii a čtečku kódů za zásadní funkce.
R4 uvádí, že využívá oblíbené a duplikaci předešlého záznamu a ocení našeptávání dříve zadaných názvů.
R1 doplňuje potřebu rychlých zkratek, které umožní udržet zápis v~nízkém počtu kroků.

\paragraph{Domácí vaření a odhad množství.}
Domácí vaření se ukazuje jako jeden z~možných důvodů, proč uživatelé zapisování opouštějí.
R3 popisuje, že vaření od oka ztěžuje odhad gramáže a že největší bariéra nastává, když položka chybí v~databázi a je nutné ji vytvořit manuálně.
R2 zdůrazňuje potřebu práce s~plánovanými ingrediencemi a následnou jednoduchou úpravou skutečně použitých množství, včetně drobných ingrediencí typu olej nebo koření.
To vede k~požadavku na funkcionalitu, která podporuje plánování a následné rychlé dorovnání skutečně použitých množství surovin.

\paragraph{Kategorie jídel a práce se zpětným zápisem.}
Respondenti se liší v~tom, zda zapisují před jídlem nebo zpětně.
R2 u~jídel z~restaurace preferuje doplnění až doma mimo sociální situaci.
R3 uvádí, že zapisovala i večer, případně plánovala dopředu, ale nedodržení plánů vedlo ke zdvojení práce.
Z~toho vyplývá potřeba podporovat jak okamžitý zápis, tak zpětný zápis z~galerie, a zároveň nevnucovat automatické zařazování kategorií podle času zadání.

\paragraph{Přehledy, motivace a dotazování nad daty.}
Základem je pro všechny denní přehled, tedy rychlá kontrola kalorií a makroživin.
R1 požaduje jednoduché přehledy po dnech, týdnech a měsících a možnost pokládat dotazy typu porovnání období podle makroživin.
R4 vnímá týdenní souhrn jako užitečné doplnění a navrhuje dotazování přirozeným jazykem s~personalizovanými doporučeními.
R2 uvádí, že motivační měsíční souhrny formou notifikace nebo e-mailu by mu pomohly udržet rutinu.
R3 klade důraz na dietní omezení a schopnost odpovědět, kolikrát došlo k~porušení diety s~proklikem na konkrétní dny.

\paragraph{Soukromí, ukládání dat a kontrola.}
Otázka soukromí se objevuje ve smyslu kontroly nad daty a transparentností práce s~nimi.
R1 požaduje jasné funkce pro mazání a informace o~tom, co se s~daty děje.
R4 je ochotna akceptovat automatické promazávání nejstarších záznamů mimo oblíbené položky.
R3 považuje za důležité mít možnost smazat jednotlivé záznamy, i když je pro ni umístění dat méně podstatné.

\paragraph{Integrace aktivity.}
Integrace se sportovními službami a hodinkami se opakuje u~všech respondentů, byť s~rozdílnou hloubkou.
R1 chce jednoduché propojení a zobrazení příjmu versus výdeje.
R2 očekává přebírání spálených kalorií s~možností nastavit, zda a jak se výdej promítne do denního limitu.
R4 zmiňuje zájem o~integraci s~Garmin a zdůrazňuje, že ji zajímá zejména promítnutí spálených kalorií do denního přehledu.

\subsection{Shrnutí kapitoly}
\label{subsec:shrnutiUS}

Uživatelská studie ukázala, že hlavní bariérou dlouhodobého zapisování je kombinace časové náročnosti, nejistoty výsledků a slabé podpory pro opakované situace, zejména opakovaná jídla, balené produkty, domácí vaření a pokrmy z~restaurací.
Rozpoznávání pomocí umělé inteligence je vnímáno jako přínosné, ale pouze za předpokladu, že uživatel má kontrolu nad výsledkem a rozumí nejistotě spojené s~vyhodnocením.
Výstupy rozhovorů byly formulovány do scénářů případů užití a do sady požadavků, které tvoří přímý podklad pro návrh a realizaci prototypu v~kapitole~\ref{ch:navrh}.
